% ============================================================
% Chapitre 10 : Introduction a la Statistique Bayesienne
% ============================================================
\chapter{Introduction \`a la Statistique Bay\'esienne}
\label{ch:bayesienne}

\begin{center}
\textit{<<~La probabilit\'e est le langage de l'incertitude.~>>} --- Pierre-Simon de Laplace
\end{center}

% ------------------------------------------------------------
\section{Exemple motivant}
% ------------------------------------------------------------

Un m\'edecin dispose d'un test de d\'epistage dont la sensibilit\'e est 95\% et la sp\'ecificit\'e 90\%. Le patient est positif. Quelle est la probabilit\'e qu'il soit vraiment malade, sachant que la pr\'evalence est de 1\% ? Le th\'eor\`eme de Bayes donne une r\'eponse surprenante : environ 8{,}7\%, bien loin des 95\% na\"ivement attendus.

% ------------------------------------------------------------
\section{Le paradigme bay\'esien}
% ------------------------------------------------------------

\begin{definition}[Approche bay\'esienne]
\index{statistique bayesienne@statistique bay\'esienne}
En statistique bay\'esienne, le param\`etre $\theta$ est trait\'e comme une \textbf{variable al\'eatoire}. On lui associe :
\begin{itemize}
\item Une \textbf{loi a priori} $\pi(\theta)$ : croyances avant observation.
\item Les \textbf{donn\'ees} : $\mathbf{x}=(x_1,\ldots,x_n)$, de vraisemblance $L(\theta;\mathbf{x})$.
\item La \textbf{loi a posteriori} : mise \`a jour des croyances.
\end{itemize}
\end{definition}

\begin{theoreme}[Th\'eor\`eme de Bayes]
\index{Bayes}
\[
\pi(\theta\mid\mathbf{x}) = \frac{L(\theta;\mathbf{x})\cdot\pi(\theta)}{\int L(\theta;\mathbf{x})\pi(\theta)\,d\theta} \propto L(\theta;\mathbf{x})\cdot\pi(\theta).
\]
\[
\boxed{\text{Post\'erieur} \propto \text{Vraisemblance} \times \text{Prieur}}
\]
\end{theoreme}

\begin{intuition}
L'approche bay\'esienne permet d'incorporer de l'information pr\'ealable et produit une distribution compl\`ete d'incertitude sur $\theta$, pas juste une estimation ponctuelle.
\end{intuition}

% ------------------------------------------------------------
\section{Choix de la loi a priori}
% ------------------------------------------------------------

\begin{definition}[Prieur conjugu\'e]
\index{prieur conjugue@prieur conjugu\'e}
La famille $\pi(\theta)$ est \textbf{conjugu\'ee} \`a la vraisemblance si la loi a posteriori appartient \`a la m\^eme famille que le prieur.
\end{definition}

\begin{center}
\begin{tabular}{llll}
\toprule
\textbf{Vraisemblance} & \textbf{Prieur conjugu\'e} & \textbf{Post\'erieur} \\
\midrule
$\mathcal{B}(1,p)$ & $\mathrm{Beta}(\alpha,\beta)$ & $\mathrm{Beta}(\alpha+\sum x_i,\;\beta+n-\sum x_i)$ \\
$\mathcal{P}(\lambda)$ & $\Gamma(\alpha,\beta)$ & $\Gamma(\alpha+\sum x_i,\;\beta+n)$ \\
$\mathcal{N}(\mu,\sigma^2)$ ($\sigma^2$ connu) & $\mathcal{N}(\mu_0,\tau^2)$ & $\mathcal{N}(\mu_n,\tau_n^2)$ \\
$\mathcal{E}(\lambda)$ & $\Gamma(\alpha,\beta)$ & $\Gamma(\alpha+n,\;\beta+\sum x_i)$ \\
\bottomrule
\end{tabular}
\end{center}

\begin{definition}[Prieur non informatif]
\index{prieur non informatif}
Un prieur <<~vague~>> qui laisse les donn\'ees dominer. Exemples :
\begin{itemize}
\item \textbf{Prieur uniforme} : $\pi(\theta)\propto 1$ (Laplace).
\item \textbf{Prieur de Jeffreys} : $\pi(\theta)\propto\sqrt{I(\theta)}$ (invariant par reparam\'etrisation).
\end{itemize}
\end{definition}

% ------------------------------------------------------------
\section{Exemples fondamentaux}
% ------------------------------------------------------------

\begin{exemple}[Mod\`ele Beta-Binomial]
$X_1,\ldots,X_n\iid\mathcal{B}(1,p)$, prieur $p\sim\mathrm{Beta}(\alpha,\beta)$.

Post\'erieur : $p\mid\mathbf{x}\sim\mathrm{Beta}(\alpha+s,\;\beta+n-s)$ o\`u $s=\sum x_i$.

Moyenne a posteriori :
\[
\E[p\mid\mathbf{x}] = \frac{\alpha+s}{\alpha+\beta+n} = \underbrace{\frac{n}{\alpha+\beta+n}}_{\text{poids donn\'ees}}\cdot\frac{s}{n} + \underbrace{\frac{\alpha+\beta}{\alpha+\beta+n}}_{\text{poids prieur}}\cdot\frac{\alpha}{\alpha+\beta}.
\]
C'est une \textbf{moyenne pond\'er\'ee} entre l'EMV $s/n$ et la moyenne a priori $\alpha/(\alpha+\beta)$.
\end{exemple}

\begin{exemple}[Mod\`ele Normal-Normal]
$X_1,\ldots,X_n\iid\mathcal{N}(\mu,\sigma^2)$ ($\sigma^2$ connu), prieur $\mu\sim\mathcal{N}(\mu_0,\tau^2)$.

Post\'erieur : $\mu\mid\mathbf{x}\sim\mathcal{N}(\mu_n,\tau_n^2)$ avec :
\[
\mu_n = \frac{\frac{\mu_0}{\tau^2}+\frac{n\bar{x}}{\sigma^2}}{\frac{1}{\tau^2}+\frac{n}{\sigma^2}}, \qquad \frac{1}{\tau_n^2}=\frac{1}{\tau^2}+\frac{n}{\sigma^2}.
\]
La pr\'ecision a posteriori = pr\'ecision a priori + pr\'ecision des donn\'ees.
\end{exemple}

% ------------------------------------------------------------
\section{Estimateurs bay\'esiens}
% ------------------------------------------------------------

\begin{definition}[Estimateurs ponctuels bay\'esiens]
\begin{itemize}
\item \textbf{Moyenne a posteriori} : $\hat{\theta}_{\mathrm{Bayes}}=\E[\theta\mid\mathbf{x}]$ (minimise le risque quadratique).
\item \textbf{M\'ediane a posteriori} : minimise le risque en valeur absolue.
\item \textbf{Mode a posteriori} (MAP) : $\hat{\theta}_{\mathrm{MAP}}=\argmax_\theta\pi(\theta\mid\mathbf{x})$.
\end{itemize}
\end{definition}

\begin{definition}[Intervalle de cr\'edibilit\'e]
\index{intervalle de credibilite@intervalle de cr\'edibilit\'e}
Un \textbf{intervalle de cr\'edibilit\'e} \`a $100(1-\alpha)\%$ est $[a,b]$ tel que $\Prob(\theta\in[a,b]\mid\mathbf{x})=1-\alpha$. L'intervalle HPD (Highest Posterior Density) est le plus court.
\end{definition}

\begin{attention}
L'intervalle de cr\'edibilit\'e a une interpr\'etation directe : <<~la probabilit\'e que $\theta\in[a,b]$ est $1-\alpha$~>>, contrairement \`a l'IC fr\'equentiste. Mais cette interpr\'etation repose sur le choix du prieur.
\end{attention}

% ------------------------------------------------------------
\section{Impl\'ementation Python}
% ------------------------------------------------------------

\begin{python}
\begin{minted}{python}
import numpy as np
from scipy import stats
import matplotlib.pyplot as plt

# --- Modele Beta-Binomial ---
# Prieur : Beta(2, 2) (leger a priori pour p ~ 0.5)
alpha_prior, beta_prior = 2, 2
n, s = 100, 68  # 68 succes sur 100

# Posterieur
alpha_post = alpha_prior + s
beta_post = beta_prior + n - s

p_grid = np.linspace(0, 1, 500)
prior = stats.beta.pdf(p_grid, alpha_prior, beta_prior)
likelihood = stats.binom.pmf(s, n, p_grid)
posterior = stats.beta.pdf(p_grid, alpha_post, beta_post)

fig, ax = plt.subplots(figsize=(10, 6))
ax.plot(p_grid, prior/prior.max(), 'b--', lw=2, label='Prieur Beta(2,2)')
ax.plot(p_grid, likelihood/likelihood.max(), 'g:', lw=2, label='Vraisemblance')
ax.plot(p_grid, posterior/posterior.max(), 'r-', lw=2, label=f'Post. Beta({alpha_post},{beta_post})')
ax.axvline(s/n, color='gray', linestyle=':', label=f'EMV = {s/n:.2f}')
ax.set_xlabel('p')
ax.set_ylabel('Densite (normalisee)')
ax.set_title('Inference bayesienne : modele Beta-Binomial')
ax.legend()
plt.savefig("ch10_beta_binomial.pdf")
plt.show()

# Estimateurs bayesiens
mean_post = alpha_post / (alpha_post + beta_post)
median_post = stats.beta.median(alpha_post, beta_post)
mode_post = (alpha_post - 1) / (alpha_post + beta_post - 2)
print(f"Moyenne posterieure : {mean_post:.4f}")
print(f"Mediane posterieure : {median_post:.4f}")
print(f"Mode (MAP)          : {mode_post:.4f}")
print(f"EMV                 : {s/n:.4f}")

# Intervalle de credibilite 95%
ci_bayes = stats.beta.interval(0.95, alpha_post, beta_post)
print(f"IC credibilite 95%  : [{ci_bayes[0]:.4f}, {ci_bayes[1]:.4f}]")

# --- Modele Normal-Normal ---
sigma2 = 4  # variance connue
mu0, tau2 = 0, 10  # prieur N(0, 10)
n_obs = 20
x_data = np.random.normal(3, np.sqrt(sigma2), n_obs)
xbar = np.mean(x_data)

prec_prior = 1/tau2
prec_data = n_obs/sigma2
prec_post = prec_prior + prec_data
mu_post = (mu0 * prec_prior + xbar * prec_data) / prec_post
tau2_post = 1 / prec_post

print(f"\nNormal-Normal:")
print(f"Moyenne posterieure : {mu_post:.4f}")
print(f"Variance posterieure: {tau2_post:.4f}")
print(f"Xbar (EMV)          : {xbar:.4f}")
\end{minted}
\end{python}

% ------------------------------------------------------------
\section{Exercices}
% ------------------------------------------------------------

\begin{exercice}[$\star$ -- Beta-Binomial]
Un joueur de basket marque 7 paniers sur 10 lancers. Avec un prieur $\mathrm{Beta}(1,1)$ (uniforme), calculer la loi a posteriori, la moyenne et un intervalle de cr\'edibilit\'e \`a 90\%.
\end{exercice}

\begin{exercice}[$\star\star$ -- Normal-Normal]
D\'eriver la loi a posteriori dans le mod\`ele Normal-Normal en compl\'etant le carr\'e dans l'exposant.
\end{exercice}

\begin{exercice}[$\star\star$ -- Prieur de Jeffreys]
Calculer le prieur de Jeffreys pour : (a) $\mathcal{B}(1,p)$, (b) $\mathcal{P}(\lambda)$, (c) $\mathcal{N}(\mu,\sigma^2)$ ($\sigma^2$ connu).
\end{exercice}

\begin{exercice}[$\star\star\star$ -- Projet]
Comparer l'approche fr\'equentiste (EMV + IC) et l'approche bay\'esienne (moyenne post\'erieure + IC de cr\'edibilit\'e) pour l'estimation de $p$ dans un mod\`ele binomial, en variant le prieur et la taille d'\'echantillon. Visualiser la convergence vers le m\^eme r\'esultat quand $n\to\infty$.
\end{exercice}

% ------------------------------------------------------------
\begin{formules}
\begin{align*}
\pi(\theta\mid\mathbf{x}) &\propto L(\theta;\mathbf{x})\cdot\pi(\theta) \\
\text{Beta-Binomial} &: \mathrm{Beta}(\alpha,\beta)\to\mathrm{Beta}(\alpha+s,\beta+n-s) \\
\text{Normal-Normal} &: \mu_n=\frac{\mu_0/\tau^2+n\bar{x}/\sigma^2}{1/\tau^2+n/\sigma^2} \\
\text{IC cr\'edibilit\'e} &: \Prob(\theta\in[a,b]\mid\mathbf{x})=1-\alpha
\end{align*}
\end{formules}
